\appendix
\chapter{Phụ lục A — Từ điển thuật ngữ}

Bảng dưới liệt kê đầy đủ thuật ngữ tiếng Anh xuất hiện trong mã nguồn và trong báo cáo,
kèm giải thích tiếng Việt.

\section{Kiến trúc và hạ tầng}

\begin{longtable}{L{4.2cm} L{9.4cm}}
\toprule
\textbf{Thuật ngữ} & \textbf{Giải thích}\\
\midrule
\endhead
\en{Monorepo} & Một kho mã duy nhất chứa nhiều dịch vụ, cho phép sửa xuyên dịch vụ trong một lần \en{commit}.\\
\en{Microservice} & Dịch vụ nhỏ, triển khai độc lập, giao tiếp qua mạng.\\
\en{API Gateway} & Cổng duy nhất nhận mọi yêu cầu từ bên ngoài rồi phân phối tới dịch vụ phía sau.\\
\en{Middleware} & Lớp bọc quanh ứng dụng, xử lý mọi yêu cầu trước và sau ứng dụng.\\
\en{Reverse proxy} & Máy chủ đứng trước ứng dụng, nhận yêu cầu thay cho ứng dụng.\\
\en{Rate limiting} & Giới hạn số yêu cầu trong một khoảng thời gian.\\
\en{Sliding window} & Cửa sổ trượt — cách tính giới hạn tần suất theo khoảng thời gian trôi liên tục.\\
\en{WAF} (\en{Web Application Firewall}) & Tường lửa ứng dụng web, lọc yêu cầu độc hại ở tầng biên.\\
\en{CORS} (\en{Cross-Origin Resource Sharing}) & Cơ chế cho phép trang web ở miền này gọi API ở miền khác.\\
\en{Preflight} & Yêu cầu \code{OPTIONS} trình duyệt gửi trước để hỏi phép trong CORS.\\
\en{Container} & Đơn vị đóng gói ứng dụng cùng phụ thuộc, chạy cách ly.\\
\en{Health check} & Điểm cuối để hệ thống điều phối biết dịch vụ còn sống và sẵn sàng.\\
\en{Liveness / Readiness} & Còn sống / sẵn sàng nhận lưu lượng — hai câu hỏi khác nhau.\\
\en{Smoke test} & Kiểm thử khói — bộ kiểm tra nhanh xem hệ thống có hoạt động ở mức cơ bản.\\
\en{Rollback} & Quay lui về phiên bản trước.\\
\en{CI/CD} & Tích hợp liên tục / triển khai liên tục.\\
\bottomrule
\end{longtable}

\section{Dữ liệu và độ tin cậy}

\begin{longtable}{L{4.2cm} L{9.4cm}}
\toprule
\textbf{Thuật ngữ} & \textbf{Giải thích}\\
\midrule
\endhead
\en{ORM} (\en{Object-Relational Mapping}) & Ánh xạ bảng cơ sở dữ liệu thành lớp trong ngôn ngữ lập trình.\\
\en{Migration} & Di trú — kịch bản thay đổi lược đồ cơ sở dữ liệu có phiên bản.\\
\en{Transaction} & Giao dịch — nhóm thao tác hoặc thành công toàn bộ hoặc không có gì xảy ra.\\
\en{UPSERT} & Chèn nếu chưa có, cập nhật nếu đã có.\\
\en{Idempotent} & Bất biến khi lặp — gọi nhiều lần cho cùng một kết quả.\\
\en{Outbox pattern} & Ghi sự kiện vào cùng giao dịch với dữ liệu, rồi mới phát đi — bảo đảm không mất sự kiện.\\
\en{Spool} & Vùng đệm ghi trước, dùng như hàng đợi bền vững.\\
\en{Lease} & Thuê — cơ chế một tiến trình giữ quyền xử lý một mục trong thời gian giới hạn.\\
\codesp{FOR UPDATE SKIP LOCKED} & Câu lệnh SQL cho phép nhiều tiến trình lấy các dòng khác nhau mà không chờ nhau.\\
\en{Epoch} & Số nguyên đánh dấu thế hệ trạng thái; tăng lên khiến mọi thứ gắn thế hệ cũ bị coi là lỗi thời.\\
\en{TTL} (\en{Time To Live}) & Thời gian sống của một mục trong bộ nhớ đệm.\\
\en{LRU} (\en{Least Recently Used}) & Chính sách loại bỏ mục lâu không dùng nhất.\\
\en{Cache invalidation} & Vô hiệu hoá bộ đệm — làm cho dữ liệu đệm cũ không còn được dùng.\\
\en{Checksum / sha256} & Tổng kiểm tra, dùng để xác minh tệp không bị đổi.\\
\en{Provenance} & Xuất xứ dữ liệu — vì sao, từ đâu và bằng cách nào một dữ liệu có mặt.\\
\bottomrule
\end{longtable}

\section{Trí tuệ nhân tạo}

\begin{longtable}{L{4.2cm} L{9.4cm}}
\toprule
\textbf{Thuật ngữ} & \textbf{Giải thích}\\
\midrule
\endhead
\en{LLM} & Mô hình ngôn ngữ lớn.\\
\en{Prompt} & Lời nhắc — văn bản đầu vào định hướng mô hình.\\
\en{System prompt} & Lời nhắc hệ thống, đặt vai trò và ràng buộc cho mô hình.\\
\en{Token} & Đơn vị văn bản mô hình xử lý; chi phí và giới hạn đều tính theo đơn vị này.\\
\en{Context window} & Cửa sổ ngữ cảnh — lượng token tối đa mô hình xử lý một lần.\\
\en{Inference} & Suy luận — chạy mô hình để sinh kết quả.\\
\en{Quantization} & Lượng tử hoá — giảm độ chính xác số để mô hình nhẹ và nhanh hơn.\\
\en{Embedding} & Véc-tơ số biểu diễn ngữ nghĩa của văn bản.\\
\en{Cosine similarity} & Độ tương đồng cô-sin giữa hai véc-tơ, đo mức giống nhau về ngữ nghĩa.\\
\en{Vector store} & Kho véc-tơ hỗ trợ tìm kiếm theo độ tương đồng.\\
\en{RAG} & Sinh có truy xuất tài liệu.\\
\en{CAG} & Sinh có bộ nhớ đệm.\\
\en{GraphRAG} & RAG kết hợp đồ thị tri thức.\\
\en{Grounding} & Nối đất — buộc câu trả lời dựa trên bằng chứng cụ thể thay vì bịa.\\
\en{Hallucination} & Ảo giác — mô hình bịa ra thông tin không có căn cứ.\\
\en{Fallback} & Đường dự phòng khi lựa chọn chính thất bại.\\
\en{Lazy loading} & Nạp lười — chỉ nạp khi thực sự cần dùng.\\
\en{Warmup} & Làm ấm — nạp trước mô hình để yêu cầu đầu tiên không bị chậm.\\
\en{TTFT} (\en{Time To First Token}) & Thời gian tới token đầu tiên — chỉ số cảm nhận độ nhanh.\\
\en{Streaming} & Phát theo dòng — trả kết quả từng mảnh thay vì đợi trọn vẹn.\\
\en{StateGraph} & Đồ thị trạng thái điều phối các bước xử lý.\\
\en{Agent} & Tác tử — chương trình dùng mô hình để tự quyết định chuỗi hành động.\\
\en{BFS} (\en{Breadth-First Search}) & Duyệt đồ thị theo chiều rộng.\\
\en{Jaccard index} & Hệ số Jaccard — kích thước giao chia kích thước hợp của hai tập hợp.\\
\en{NDCG, MRR, Precision@k, Recall@k} & Bốn chỉ số đánh giá chất lượng xếp hạng kết quả truy xuất.\\
\en{WER} (\en{Word Error Rate}) & Tỉ lệ lỗi từ.\\
\en{TTR} (\en{Type-Token Ratio}) & Tỉ lệ từ khác nhau trên tổng số từ, đo độ phong phú từ vựng.\\
\bottomrule
\end{longtable}

\section{Sư phạm và mô hình người học}

\begin{longtable}{L{4.2cm} L{9.4cm}}
\toprule
\textbf{Thuật ngữ} & \textbf{Giải thích}\\
\midrule
\endhead
\en{CEFR} & Khung tham chiếu ngôn ngữ chung châu Âu, sáu bậc A1--C2.\\
\en{Spaced repetition} & Lặp lại ngắt quãng.\\
\en{Forgetting curve} & Đường quên — mô tả xác suất nhớ giảm theo thời gian.\\
\en{Spacing effect} & Hiệu ứng giãn cách — ôn cách quãng hiệu quả hơn ôn dồn.\\
\en{FSRS} & Bộ lập lịch lặp lại ngắt quãng tự do, mô hình hoá trí nhớ qua độ bền và độ khó.\\
\en{SM-2} & Thuật toán lặp lại ngắt quãng của SuperMemo, thế hệ trước FSRS.\\
\en{Stability} & Độ bền trí nhớ tính bằng ngày.\\
\en{Retrievability} & Xác suất nhớ lại tại thời điểm hiện tại.\\
\en{Mastery} & Mức nắm vững một khái niệm.\\
\en{BKT} (\en{Bayesian Knowledge Tracing}) & Mô hình theo dõi tri thức kiểu Bayes, có tham số trượt và đoán mò — \textbf{không} phải thứ hệ thống này dùng.\\
\en{Delta rule} & Luật delta — cập nhật giá trị theo sai lệch giữa mục tiêu và giá trị hiện tại.\\
\en{Slip / Guess} & Trượt (biết nhưng làm sai) / đoán mò (không biết nhưng làm đúng).\\
\en{Scaffolding} & Giàn giáo — chiến lược dạy chia nhỏ hỗ trợ rồi rút dần.\\
\en{Socratic} & Kiểu hỏi gợi mở dẫn người học tự tìm ra câu trả lời.\\
\en{Prerequisite} & Kiến thức tiên quyết.\\
\en{Gamification} & Trò chơi hoá — dùng cơ chế trò chơi để tạo động lực học.\\
\en{Streak} & Chuỗi ngày học liên tiếp.\\
\en{XP} (\en{Experience Points}) & Điểm kinh nghiệm.\\
\en{Leaderboard} & Bảng xếp hạng.\\
\en{Grinding} & Cày điểm — lặp hành động dễ để tăng chỉ số mà không thực sự học.\\
\bottomrule
\end{longtable}

\chapter{Phụ lục B — Bảng tra nhanh}

\section{Hằng số thuật toán}

\begin{longtable}{L{5.4cm} L{2.6cm} L{5.6cm}}
\toprule
\textbf{Hằng số} & \textbf{Giá trị} & \textbf{Nơi dùng}\\
\midrule
\endhead
\code{ALGORITHM\_VERSION} & \code{delta-fsrs-v3} & Trạng thái khái niệm\\
\code{LEARNING\_RATE} \(\eta\) & 0,45 & Luật delta\\
\code{FORGETTING\_DECAY} \(D\) & \(-0{,}5\) & Đường quên\\
\code{FORGETTING\_FACTOR} \(F\) & \(19/81\) & Đường quên\\
\code{TARGET\_RETENTION} & 0,9 & Khoảng ôn tập\\
\code{STABILITY\_GROWTH} & 0,45 & Tăng độ bền khi đúng\\
\code{SPACING\_BONUS} & 1,6 & Thưởng giãn cách\\
\code{LAPSE\_DECAY} & 0,55 & Giảm độ bền khi sai\\
\code{DIFFICULTY\_EASE\_STEP} & 0,04 & Giảm độ khó khi đúng\\
\code{DIFFICULTY\_LAPSE\_STEP} & 0,08 & Tăng độ khó khi sai\\
\code{MIN/MAX\_MASTERY} & 0,01 / 0,99 & Kẹp mức nắm vững\\
\code{MIN/MAX\_STABILITY\_DAYS} & 0,25 / 3650 & Kẹp độ bền\\
\code{CONFIDENCE\_FULL\_EXERCISES} & 50 & Độ tin cậy điểm kỹ năng\\
\code{MAX\_SCORE\_STEP} & 0,15 & Trần một bước EMA\\
\code{decay\_factor} \(\lambda\) & 0,95 & Trọng số lịch sử EMA\\
\code{\_TAU\_REUSE} \(\tau_0\) & 0,25 & Ngưỡng tái dùng đệm\\
\code{\_TAU\_PATCH} \(\tau_1\) & 0,55 & Ngưỡng vá đệm\\
\code{concept\_floor} & 0,50 & Sàn chồng lấp khái niệm\\
\code{\_FUSION\_ALPHA/BETA/GAMMA} & 0,5 / 0,3 / 0,2 & Hoà trộn truy xuất\\
\code{\_RECENCY\_LAMBDA} & 0,01 & Suy giảm độ mới\\
\code{DAILY\_XP\_CAP} & 500 & Trần XP ngày\\
\code{MAX\_SINGLE\_AWARD} & 200 & Trần một lần cấp XP\\
\bottomrule
\caption{Hằng số thuật toán và nơi sử dụng}
\end{longtable}

\section{Tệp cần biết}

\begin{longtable}{L{7.6cm} L{6.0cm}}
\toprule
\textbf{Đường dẫn} & \textbf{Vai trò}\\
\midrule
\endhead
\code{backend-service/app/services/learner\_state.py} & Động cơ FSRS + luật delta\\
\code{backend-service/app/services/proficiency\_service.py} & Chấm điểm CEFR\\
\code{backend-service/app/services/xp\_service.py} & XP và chống cày điểm\\
\code{backend-service/app/services/content\_agent\_apply.py} & Đường duy nhất ghi khoá học\\
\code{backend-service/app/routes/admin\_courses.py} & \code{\_apply\_lesson\_update}\\
\code{backend-service/scripts/audit\_course\_content.py} & Kiểm toán và sửa nội dung\\
\code{backend-service/scripts/generate\_exercises\_ai.py} & Sinh và sinh lại bài tập\\
\code{ai-service/api/services/trace\_cag/graph.py} & Đồ thị pipeline\\
\code{ai-service/api/services/trace\_cag/cache\_utils.py} & Cổng đệm RAPID\\
\code{ai-service/api/services/trace\_cag/l1\_state\_cache.py} & Quyết định SCAR-L1\\
\code{ai-service/api/services/trace\_cag/retrieve.py} & Hoà trộn truy xuất\\
\code{ai-service/api/services/content\_etl/} & ETL ngữ liệu\\
\code{flutter-app/lib/features/learning/presentation/screens/learning\_roadmap\_screen.dart} & Lộ trình học\\
\code{flutter-app/lib/core/services/skill\_event\_recorder.dart} & Gửi sự kiện học tập\\
\code{scripts/deploy-one-shot.sh} & Triển khai chuẩn\\
\code{scripts/restore-prod.sh} & Khôi phục có kiểm tra trước\\
\bottomrule
\caption{Tệp then chốt trong kho mã}
\end{longtable}
